\subsection{积空间}

定义 1. 给定拓扑空间的一个族 \((\mathbf{X}_{t})_{t\in I}\)，这个族的积空间是指积集合 \(\mathbf{X} = \prod_{t\in I}\mathbf{X}_{t}\) 赋以诸 \(\mathbf{X}_{t}\) 的拓扑的积（§ 2，第 3 目，例 3）。诸空间 \(\mathbf{X}_{t}(t\in I)\) 称为 \(\mathbf{X}\) 的因子。

由 § 2 第 3 目命题 4，X 上的积拓扑以形如 \(\overline{\mathrm{pr}}_{t}^{1}(\mathbf{U}_{t})\) 的集合的有限交组成的集合 B 为一个基，其中 \(U_{t}\) 在 \(X_{t}\) 中是开集；这些集合是积 \(\prod_{i\in I}A_{i}\)，其中 \(A_{i}\) 在 \(X_{t}\) 中是开集（对每个 \(i\in I\)），且除去有限多个指标外 \(A_{i}=X_{t}\)。这些集合将称为初等集。

若 \(\mathfrak{B}_i\) 是 \(X_i\) 的拓扑的一个基（对每个 \(i \in I\)），则显然，初等集 \(\prod_{i \in I} A_i\)（满足 \(A_i \in \mathfrak{B}_i\)（对每个指标 \(i\)，只要 \(A_i \neq X_i\)））组成积拓扑的另一个基。于是，包含给定点 \(x \in X\) 的这种类型的初等集组成 \(x\) 的一个基本邻域系（\(\S 1\)，第 3 目，命题 3）。

若 I 是有限集，由诸因子 \(X_{i}\) 的拓扑构造积拓扑就更简单：初等集恰是积 \(\prod_{i\in I}A_{i}\)，其中 \(A_{i}\) 是 \(X_{i}\) 的任意开子集（对每个 \(i\in I\)）（参见习题 9）。

例。* 积 \(\mathbf{R}^n\)（\(n\) 个与实直线 \(\mathbf{R}\) 恒同的空间的积）称为 \(n\) 维实数空间；\(\mathbf{R}^2\) 也称为实平面（参见第六章，§ 1）。类似地，从有理直线 \(\mathbf{Q}\) 出发，我们定义 \(n\) 维有理数空间 \(\mathbf{Q}^n\)（当 \(n = 2\) 时称为有理平面）。

空间 \(\mathbf{R}^n\) 的拓扑以 \(n\) 个 \(\mathbf{R}\) 中开区间的所有积组成的集合为一个基，这些积称为 \(n\) 维开方块。包含点 \(x \in \mathbf{R}^n\) 的开方块组成这一点的一个基本邻域系。类似地，\(n\) 个 \(\mathbf{R}\) 中闭区间的积称为 \(n\) 维闭方块。以 \(x\) 为内点的闭方块也组成 \(x\) 的一个基本邻域系。对 \(\mathbf{Q}^n\) 有类似的结果。*

命题 1. 设 \(f = (f_{i})\) 是拓扑空间 \(Y\) 到积空间 \(X = \prod_{i \in I} X_{i}\) 中的一个映射。则 \(f\) 在点 \(a \in Y\) 处连续，当且仅当 \(f_{i}\)（对每个 \(i \in I\)）在 a 处连续。

由于 \(f_{t} = \mathrm{pr}_{t} \circ f\)，这只是 §2 第 3 目命题 4 的一个特例。

推论 1. 设 \((\mathbf{X}_t)_{i \in \mathbf{I}}\)、\((\mathbf{Y}_t)_{i \in \mathbf{I}}\) 是具有同一指标集的两个拓扑空间族。对每个 \(i \in \mathbf{I}\)，设 \(f_i\) 是 \(\mathbf{X}_t\) 到 \(\mathbf{Y}_t\) 中的一个映射。为使积映射 \(f: (x_i) \to (f_i(x_i))\)（它把

\[
\prod_ {i \in I} X _ {i} \quad i n t o \quad \prod_ {i \in I} Y _ {i}
\]

映到）在一点 \(a = (a_{i})\) 处连续，必须且只需 \(f_{i}\) 在 \(a_{i}\) 处连续（对每个 \(\mathfrak{t} \in \mathbf{I}\)）。

\(f\) 可以写成 \(x \to (f_{t}(\mathrm{pr}_{t}(x))\)，于是由命题 1，该条件是充分的。反之，对每个 \(x \in I\)，设 \(g_{x}\) 是 \(X_{x}\) 到 \(\prod_{i \in I} X_{i}\) 中的这样的映射：\(\mathrm{pr}_{x}(g_{x}(x_{x})) = x_{x}\)，且 \(\mathrm{pr}_{t}(g_{x}(x_{x})) = a_{t}\)（每当 \(i \neq x\)）；则由命题 1，\(g_{x}\) 在点 \(a_{x}\) 处连续。由于 \(f_{x} = \mathrm{pr}_{x} \circ f \circ g_{x}\)，由此得出：若 \(f\) 在 \(a\) 处连续，则 \(f_{x}\) 在 \(a_{x}\) 处连续。

推论 2. 设 X、Y 是两个拓扑空间。为使映射 \(f: \mathbf{X} \to \mathbf{Y}\) 连续，必须且只需映射 \(g: x \to (x, f(x))\) 是 X 到 \(f\) 的图象 G 上（把 G 看作积空间 \(\mathbf{X} \times \mathbf{Y}\) 的子空间）的一个同胚。

由于 \(f = \mathrm{pr}_2 \circ g\)，该条件是充分的。它也是必要的，因为若 \(f\) 连续，则 \(g\) 是双射且连续（命题 1），且 \(g\) 的逆是 \(\mathrm{pr}_1\) 在 \(G\) 上的限制，它是连续的（参见《集合论》第四章，§ 2，第 4 目，判别准则 CST 17）。

命题 2（拓扑积的结合性）. 设 \((\mathbf{X}_{\iota})_{\iota \in \mathbf{I}}\) 是拓扑空间的一个族，\((\mathbf{J}_{\kappa})_{\kappa \in \mathbf{K}}\) 是集合 I 的一个分拆，且对每个 \(\kappa \in \mathbf{K}\)，设 \(\mathbf{X}_{\kappa}^{\prime} = \prod_{i\in \mathbf{J}_{\kappa}}\mathbf{X}_{\iota}\) 是诸空间 \(\mathbf{X}_{\iota}\)（对 \(\iota \in \mathbf{J}_{\mathbf{K}}\)）的积。则积空间的典范映射
% semantic-fragments: legacy-input-seam
\relax{}
\[\prod_{i \in I} X_i\ \text{到积空间}\ \prod_{x \in K} X_x'\]

是一个同胚。

这是初始拓扑的传递性的一个特殊情形（§ 2，第 3 目，命题 5；参见《集合论》，第四章，§ 2，第 4 目，判别准则 CST 13）。

通常我们借助典范映射把积空间 \(\prod_{i\in I}X_i\) 与 \(\prod_{x\in X}X_x'\) 等同起来。

推论. 设 \(\sigma\) 是集合 I 的一个置换。则映射 \((x_{i})\to (x_{\sigma(i)})\) 是

\[
\prod_ {i \in I} X _ {i} \quad onto \quad \prod_ {i \in I} X _ {\sigma (i)}.
\]

在命题 2 中取 \(\mathbf{K} = \mathbf{I}\) 及 \(\mathbf{J}_t = \{\sigma(t)\}\)（对每个 \(t \in I\)）。

命题 3. 设 X 是一个集合，\((\mathbf{Y}_{t})_{t\in I}\) 是拓扑空间的一个族，且对每个 \(t\in I\)，设 \(f_{t}\) 是 X 到 \(Y_{t}\) 中的一个映射。设 f 是映射 \(x\to(f_{t}(x))\)，把 X 映入 \(Y=\prod_{t\in I}Y_{t}\)，并设 G 是 X 上使诸映射 \(f_{t}\) 连续的最粗的拓扑。则 G 是 Y 上的积拓扑在 \(f(\mathbf{X})\) 上诱导的拓扑在 f 下的逆像。

这是初始拓扑的传递性的又一个特殊情形（§ 2，第 3 目，命题 5；参见《集合论》，第四章，§ 2，第 4 目，判别准则 CST 15）。

推论. 对每个 \(t \in I\)，设 \(A_{t}\) 是 \(Y_{t}\) 的一个子空间。则在 \(A = \prod_{t \in I} A_{t}\) 上由 \(\prod_{t \in I} Y_{t}\) 上的积拓扑诱导的拓扑是诸子空间 \(A_{t}\) 的拓扑之积。

设 \(j_{i}\) 是典范单射 \(A_{i} \to Y_{i}\)，并对诸映射 \(f_{i} = j_{i} \circ \operatorname{pr}_{i}\) 应用命题 3（参见《集合论》，第四章，§ 2，第 4 目，判别准则 GST 14）。

